\section{Supplementary Theoretical Details}

\subsection{Proof of Proposition~\ref{prop:specialist_bottleneck}}
\label{apx:specialist_bottleneck_proof}

Let $\widetilde L_c(\pi)$ denote the normalized condition-specific loss introduced
before the proposition. For $s\in\mathcal{S}$, a fixed policy $\pi_s$ chooses the
same optional channel $s$ for every operating condition. The policy
$\pi_{\text{label}}$ has privileged access to the condition label and changes to the
corresponding optional channel whenever the common startup and minimum
dwell-time constraints permit. Condition intervals are long
enough for these changes, and the losses incurred during transitions are
included in $\widetilde L_c$.

For a fixed optional channel $s$, define
$\delta_c(s)=\widetilde L_c(\pi_s)-
\widetilde L_c(\pi_{\text{label}})$ and the mismatched-condition set
\begin{equation}
  \mathcal{C}_{\mathrm{mis}}(s)
  =
  \{c\in\mathcal{C}\mid s\text{ mismatches the best optional channel in }c\}.
\end{equation}
By assumption, $\delta_c(s)\geq0$ for every condition and
$\delta_c(s)\geq\Delta_c>0$ for each mismatch. The weighted normalized forecast
loss of the fixed channel therefore satisfies
\begin{equation}
  L_{\rho}(\pi_s)
  -
  L_{\rho}(\pi_{\text{label}})
  =
  \sum_{c\in\mathcal{C}}\rho_c\delta_c(s)
  \geq
  \sum_{c\in\mathcal{C}_{\mathrm{mis}}(s)}
  \rho_c \Delta_c
  >
  0 .
\end{equation}
The strict inequality follows because
$\mathcal{C}_{\mathrm{mis}}(s)$ is non-empty and every included term is
positive. Thus every fixed channel has higher weighted normalized forecast loss under the stated
sufficient conditions.

\subsection{Proof of Proposition~\ref{prop:non_equivalence}}
\label{apx:non_equivalence_proof}

Set the common objective horizon to $T=1$; then all three objectives assign the
same temporal weight $\gamma^0=1$ to the evaluated step. Consider one forecast
target, one forecast horizon, and two independent scalar
Gaussian random walks
\begin{equation}
  z_{t+1}=z_t+\xi_t,
  \qquad
  e_{t+1}=e_t+\zeta_t,
  \label{eq:proxy_counterexample_state}
\end{equation}
where $\xi_t\sim\mathcal{N}(0,q_z)$ and
$\zeta_t\sim\mathcal{N}(0,q_e)$ are mutually independent over time. The
one-step forecast target is
\begin{equation}
  y_{t+1}=c z_t+M e_t+\nu_{t+1},
  \label{eq:proxy_counterexample_target}
\end{equation}
with independent $\nu_{t+1}\sim\mathcal{N}(0,\sigma_\nu^2)$ and coefficients
$M>c>0$. Suppose that one random-walk component can be measured without noise at
each time step. At the evaluated time step, feasible schedules can produce age vectors
$d^{(1)}=(0,1)$ under policy $\pi_1$ and $d^{(2)}=(2,0)$ under policy $\pi_2$.
Here $\mathcal{Y}=\{y\}$ for the forecast-loss objective and
$\mathcal{V}_{\mathrm{AoI}}=\{z,e\}$ for the AoI objective.
These vectors can arise from feasible one-measurement schedules that prioritize
different random-walk components before the evaluated time step.
Thus $\pi_1$ has lower total AoI because $1<2$.

The corresponding posterior error covariances are
$P^{(1)}=\operatorname{diag}(0,q_e)$ and
$P^{(2)}=\operatorname{diag}(2q_z,0)$. Choose $q_z>0$ and
\begin{equation}
  \frac{2c^2q_z}{M^2}<q_e<2q_z,
  \label{eq:proxy_counterexample_condition}
\end{equation}
which is possible because $M>c$. The right inequality gives
$\operatorname{tr}P^{(1)}<\operatorname{tr}P^{(2)}$, so the covariance trace
objective also prefers $\pi_1$.

Now take the single-target, single-horizon, positive-weight special case of
\Cref{eq:step_forecast_loss}. Its positive weight cancels between the numerator
and denominator, so the per-step forecast loss is
\begin{equation}
  L_{\mathrm{step}}(t)
  =
  \min\!\left\{10,\,
  \frac{|\widehat y_{t+1}-y_{t+1}|}{\sigma}\right\},
  \qquad \sigma>0.
  \label{eq:proxy_counterexample_step_loss}
\end{equation}
The conditional forecast errors under $\pi_1$ and $\pi_2$ are zero-mean Gaussian
with variances
\begin{equation}
  v_1=M^2q_e+\sigma_\nu^2,
  \qquad
  v_2=2c^2q_z+\sigma_\nu^2.
  \label{eq:proxy_counterexample_error_variances}
\end{equation}
The left inequality in \Cref{eq:proxy_counterexample_condition} gives
$v_1>v_2$.

For $v>0$, let $Z_v\sim\mathcal{N}(0,v)$ and define
\begin{equation}
  g(v)=\mathbb{E}\!\left[
  \min\!\left\{10,\frac{|Z_v|}{\sigma}\right\}
  \right].
\end{equation}
Writing $Z_v=\sqrt{v}\,U$ with $U\sim\mathcal{N}(0,1)$ shows that the integrand
for $v_1$ is at least that for $v_2$ at every $U$. The inequality is strict on
\begin{equation}
  0<|U|<\frac{10\sigma}{\sqrt{v_1}},
\end{equation}
which has positive probability. Hence $g(v_1)>g(v_2)$, and
\begin{equation}
  \mathbb{E}[L_{\mathrm{step}}(t;\pi_1)]
  >
  \mathbb{E}[L_{\mathrm{step}}(t;\pi_2)].
\end{equation}
Because $T=1$ in \Cref{eq:forecast_objective}, this gives
$J_{\mathrm{fcst}}(\pi_2)>J_{\mathrm{fcst}}(\pi_1)$. The AoI and covariance
trace objectives therefore prefer $\pi_1$, while the capped, scale-normalized
forecast-loss objective prefers $\pi_2$.

\subsection{Observation-dependent scheduling and fixed cycles}

A fixed cyclic schedule chooses optional channels from the time index alone. Such a schedule may be
optimal in structured settings, including periodic conditions or rewards that are
phase-aligned with the cycle. Without that structure, a fixed cycle need not be
optimal. If informative observations predict a change in operating condition and a cycle has
nonzero probability of assigning a mismatched optional channel, the positive mismatch
loss condition in Proposition~\ref{prop:specialist_bottleneck} yields positive
expected regret relative to a policy conditioned on those observations. Action
sequences are therefore needed to distinguish observation-dependent scheduling
from a fixed schedule or a fixed cyclic schedule.

\subsection{Relation to fixed channel-subset allocation}

\Cref{prop:non_equivalence} does not imply that an adaptive policy must beat a
fixed channel subset. If one fixed optional channel measures the variables that dominate a
forecast metric, a fixed policy can remain optimal for that metric. Adaptive
value requires the incompatible-condition requirement in
\Cref{prop:specialist_bottleneck}, together with observations that provide useful
information about a change in operating condition. Mean forecast loss and
weighted normalized forecast loss can therefore rank a fixed and an adaptive
policy differently.

\section{Simulation Validation}
\label{apx:generator_validation}

\Cref{tab:generator_validation} reports the acceptance checks used to bound the
simulation domain. Supplementary Fig. S3 provides the corresponding
distribution, autocorrelation, and event-window plots. These checks document
the simulation construction and are not policy-selection criteria.

\begin{table}[htbp]
  \centering
  \caption{Simulation validation checks and acceptance criteria.}
  \label{tab:generator_validation}
  \small
  \setlength{\tabcolsep}{3.5pt}
\begin{tabular}{@{}>{\raggedright\arraybackslash}p{0.43\linewidth}>{\raggedright\arraybackslash}p{0.17\linewidth}>{\raggedright\arraybackslash}p{0.24\linewidth}c@{}}
\toprule
Validation criterion & Value & Acceptance rule & Passed \\
\midrule
Wind-speed autocorrelation max deviation, lags 1--12 h & $0.048$ & $<0.05$ & yes \\
$P(\mathrm{event\ fraction}>0.75)$ in 512-step windows & $0.065$ [$0.062,0.067$] & $>0.05$ & yes \\
$P(\mathrm{event\ fraction}<0.25)$ in 512-step windows & $0.539$ & $>0.30$ & yes \\
Max KS statistic against reference distributions from Antarctic AWS stations & $0.0146$ & $<0.05$ & yes \\
Median blowing-snow event duration & $18.0$ h & $12$--$20$ h & yes \\
Flux--wind relationship, log--log slope & $2.966$ & $2.5$--$3.5$ & yes \\
Particle-size--wind correlation, Spearman & $-0.544$ & $<-0.30$ & yes \\
Wind-speed power spectrum log-MSE, 0.1--4.0 cpd & $0.0419$ & $<0.10$ & yes \\
\bottomrule
\end{tabular}

\end{table}

\clearpage
